\bookchapter{第十章：モバイル決済と生活サービス}

\chapterintro{一、銀行カードの登録方法とは？}

\booksection{AlipayでもWeChatでも、初回のカード登録は以下の手順を踏みます：}

Appを開き、「マイ」または「ウォレット」を見つける。「銀行カードを追加」をタップ。

\begin{enumerate}
  \item カード番号を入力。最近のAppはほとんど写真でカード番号を認識できるので、一桁ずつ入力する必要はありません。
  \item 銀行から事前登録した携帯番号に認証コードSMSが送信されます。入力して、6桁の数字の支払いパスワードを設定。
  \item この支払いパスワードは、銀行カードのパスワードやロック画面のパスワードとは別々にすること。楽だからといって同じにしないでください。
  \item 一つの小技：セキュリティが心配なら、銀行AppでII類またはIII類の電子口座を開設し、スマホ決済専用にしましょう。このカードには日常の小銭だけを入れておけば、問題が起きても損害は限定的です。
\end{enumerate}

二、支払いコードは「見た者が受け取る」

\booksection{多くの人が支払いコードと收款コード（受取コード）の違いを混同していますが、この違いは重要です。}

支払いコード

\begin{itemize}
  \item \textbf{：あなたがレジ係に見せてスキャンさせるバーコード＋QRコード。このコードはあなたの「電子財布」に相当し、スキャンした者はあなたの口座からお金を引き落とせます。だから}絶対にスクリーンショットを撮って誰かに送らないでください\textbf{。このコードは1分ごとに自動更新され、盗用を防いでいます。}收款コード
  \item \textbf{：他人がこのコードをスキャンしてあなたに送金します。このコードはスクショして保存でき、他人に送っても安全です。}この違いをしっかり認識すれば、多くの詐欺の罠を避けられます。「支払いコードのスクショを送って」と言う人は、100\%詐欺師です。
\end{itemize}

三、これらのセキュリティ設定、5分で済ませて安心を得る

\booksection{指紋/顔認証決済}

\begin{itemize}
  \item \textbf{：AlipayとWeChatの決済設定で、生体認証決済をオンに。これで毎回の決済時に指紋または顔スキャンが必要になり、スマホをなくしても拾った人はあなたのお金を使えません。}セキュリティロック
  \item \textbf{：AlipayとWeChatには「セキュリティロック」機能があります。オンにすると、Appの「ウォレット」ページに入る時に追加の本人確認が必要。パス：Alipay マイ 設定 アカウントとセキュリティ セキュリティセンター セキュリティロック。}暗証不要決済の確認
  \item \textbf{：多くの人が各種暗証不要決済に契約していて、自分でも忘れています。AlipayとWeChatの「決済設定」で「暗証不要決済/自動引き落とし」を探し、契約中のサービスを確認し、不要なものはオフに。}ログインデバイス管理
  \item \textbf{：定期的にどのデバイスがあなたのアカウントにログインしたか確認。知らないデバイスがあれば、すぐにキックアウト。パス：Alipay 設定 アカウントとセキュリティ ログインデバイス管理。}四、バス·地下鉄：スマホが交通カード
\end{itemize}

\booksection{最近外出してバスや地下鉄に乗る時、もう物理の交通カードを使う人は稀です。スマホを改札機にかざす方法は2つあります：}

スマホウォレットの仮想交通カード

\begin{itemize}
  \item \textbf{：Appleウォレット、Huaweiウォレット、Xiaomiウォレットで仮想交通カードを直接発行。この方法の利点は、かざす時にスマホをロック解除する必要がなく、直接近づけるだけで済むこと。応答速度が最も速く、ネット接続がなくてもかざせます。一部の都市の仮想交通カードには物理カードと同じ乗り継ぎ割引もあります。}Alipay/WeChatの乗車コード
  \item \textbf{：Appで「乗車コード」を検索し、お住まいの都市のバスや地下鉄の乗車コードを開通。この方法の利点は、チャージ不要で、残高や連携した銀行カードから引き落とされること。ただしネット接続が必要で、かざす前にAppを開いてコードを探す必要があります。}私は個人的にスマホウォレットの仮想交通カードをおすすめします。初回は10〜20元チャージする必要がありますが、使い勝手が本当に段違いに良いです。特に地下鉄に急ぐ時に、立ち止まってスマホをロック解除してAppを探す必要がありません。
\end{itemize}

五、生活費支払い、デリバリー、タクシー

\booksection{水道光熱費}

\begin{itemize}
  \item \textbf{：AlipayまたはWeChatの「生活支払い」で、ご家庭の水道、電気、ガスの戸号を連携。毎月請求書が届いたら、そこから直接支払えます。戸号はご家庭の紙の請求書で確認できるか、カスタマーサービスに電話して聞けます。}デリバリー比較
  \item \textbf{：同じ店でも美団と餓了麼で価格と割引が異なることがよくあります。注文前に両方見比べると、数元節約できることがあります。また、美団の「拼好飯」機能は、友人と一緒に注文すると価格がかなり安くなります。}タクシーの安全
  \item \textbf{：滴滴や高徳タクシーで、「緊急連絡先」と「行程共有」を設定しておきましょう。タクシーに乗ると、Appが自動的にあなたの行程を設定した連絡先に共有します。また、深夜のタクシーは「快車」より「タクシー」を優先しましょう。タクシーは業界の監督があり、安全度が高めです。}六、スマホで医療保険、社会保険、行政手続き
\end{itemize}

\booksection{ここ数年の変化は大きく、以前は窓口に並んだ多くの手続きが、今はスマホで数分で完了します。}

医療保険電子証憑

\begin{itemize}
  \item \textbf{：AlipayまたはWeChatで「医療保険電子証憑」を検索し、有効化すれば、病院受診や薬局購入時にスマホをかざすだけで済み、物理の医療保険カードを持参する必要がありません。初回は顔認証が必要で、以降は使用時に動的コードが生成されます。}社会保険照会
  \item \textbf{：「掌上12333」Appをダウンロードするか、Alipayの「市民センター」で、社会保険納付記録、年金口座残高、失業給付申請などを照会できます。}所得税
  \item \textbf{：毎年3月から6月の所得税年度調整は、「個人所得税」Appで完結。控除対象の特別追加控除（子女教育、住宅ローン利子、老人扶養など）があるか確認しましょう。これらはすべて納税額を減らしたり還付を受けたりするのに役立ちます。}これらを覚えて、詐欺を防ぎ安心を
\end{itemize}

\begin{tcbwarn}{警察、銀行、カスタマーサポートは「安全口座」への送金を指示しません。}
\begin{itemize}
  \item 誰かに支払いコードのスクショを送らせる人は、すべて詐欺師です。
  \item 見知らぬリンクをタップしない、特に「ETC期限切れ」「社会保険カード停止」などのSMS。
  \item ネットショッピングの返金は元の経路で戻り、銀行カード番号の提供は不要です。
  \item 今すぐやる：決済Appにセキュリティチェックを実施
\end{itemize}

\end{tcbwarn}

\begin{tcbpractice}{AlipayとWeChatを開き、指紋/顔認証決済が有効か確認。}
\begin{enumerate}
  \item 「セキュリティロック」機能を見つけてオンに。
  \item 「暗証不要決済/自動引き落とし」リストを確認し、不要なサービスをオフに。
  \item 「ログインデバイス」リストを確認し、知らないデバイスを削除。
  \item 医療保険電子証憑を有効化し、次回病院では物理カードを持参不要に。
  \item 第十一章：撮影と写真管理
\end{enumerate}

\end{tcbpractice}
